\section{Gap Penelitian dan Usulan Metode}

Berdasarkan penelitian terdahulu, metode IndoBERT telah terbukti
memiliki performa yang baik dalam klasifikasi sentimen berbahasa
Indonesia \cite{yulianti2023,latifah2024,wafda2025}.
Sementara itu, BERTopic mampu menghasilkan topik yang lebih koheren
dibandingkan metode tradisional seperti LDA dan NMF
\cite{babalola2023,janssens2025}.

Beberapa penelitian telah mengintegrasikan IndoBERT dan BERTopic
untuk melakukan ABSA pada domain pariwisata
\cite{kurniawan2023,girsang2024,maylawati2024}.
Namun, sebagian besar penelitian tersebut berfokus pada wisata
secara umum dan belum mengkaji wisata alam dan religi di Kabupaten
Bangkalan.

Selain itu, penelitian terdahulu umumnya menggunakan data dari satu
objek wisata atau satu kategori wisata tertentu sehingga belum mampu
memberikan gambaran yang lebih luas mengenai persepsi wisatawan
terhadap berbagai jenis destinasi wisata dalam satu wilayah.

Oleh karena itu, penelitian ini mengusulkan integrasi IndoBERT dan
BERTopic pada data Google Reviews yang diperoleh dari beberapa
destinasi wisata alam dan religi di Kabupaten Bangkalan.
IndoBERT dipilih karena kemampuannya dalam memahami konteks bahasa
Indonesia secara lebih baik dibandingkan metode konvensional
\cite{yulianti2023,latifah2024}. Sementara itu, BERTopic dipilih
karena mampu menemukan topik secara otomatis dengan kualitas topik
yang lebih baik dibandingkan LDA dan NMF
\cite{babalola2023,janssens2025}. Kombinasi kedua metode tersebut
diharapkan mampu menghasilkan pemetaan sentimen berbasis aspek yang
lebih komprehensif untuk mendukung evaluasi dan pengembangan
destinasi wisata.